\documentclass[11pt,a4paper]{article}

\usepackage[utf8]{inputenc}
\usepackage[T1]{fontenc}
\usepackage[english]{babel}
\usepackage[a4paper, margin=2.4cm]{geometry}
\usepackage{amsmath,amssymb}
\usepackage{graphicx}
\usepackage{booktabs}
\usepackage{caption}
\usepackage{subcaption}
% euro symbol rendered via UTF-8 input directly (no eurosym needed)
\usepackage{xcolor}
\usepackage{hyperref}

\hypersetup{colorlinks=true, linkcolor=blue!50!black, urlcolor=blue!50!black,
            citecolor=blue!50!black}

\title{Performance Modeling of a Cafeteria System \\
       \large A discrete-event simulation study}
\author{Daniel \\ Performance Modeling of Computer Systems and Networks \\
        Universit\`a degli Studi di Roma Tor Vergata}
\date{A.A. 2025/2026}

\begin{document}
\maketitle

\begin{abstract}
This work presents a complete simulation study of a small cafeteria
modelled as a queueing system with two coupled resources: $C$ waiters
and $S$ seats. Two customer classes are considered (coffee and
breakfast). The study follows Algorithm~1.1 of Leemis \&
Park\,\cite{leemis_park} for model development, employs the
multi-stream Lehmer pseudo-random generator, identifies the warm-up
period by Welch's method, computes all statistics via
replication-deletion with $n=32$ independent replications per scenario,
and reports $95\%$ confidence intervals based on the t-Student
distribution. The full-factorial experiment varies
$(C,S) \in \{2,3\} \times \{4,5\}$ and the arrival rate
$\lambda \in \{15,20,25,30\}$ customers/hour. Results show that
\emph{seats are the bottleneck of the system}: at the baseline rate
$\lambda=20$/h the seat utilization is roughly twice the waiter
utilization. Consequently, adding a single seat is an order of
magnitude more effective and twelve times cheaper than adding a single
waiter. An economic analysis based on operating costs and per-customer
revenue translates these findings into a clear operational policy.
\end{abstract}

\tableofcontents
\newpage

% =====================================================================
\section{Introduction and motivation}
% =====================================================================

Small service establishments such as cafeterias, fast-food counters
and bars share a common operational dilemma: when demand grows, should
the manager invest in additional service staff, in additional seating
capacity, or in both? The intuitive answer ``add a waiter when the
queue grows'' is not always correct, because the bottleneck of the
system is not necessarily the service capacity of the staff: a customer
also occupies a seat throughout the entire stay, including the time
spent consuming the product.

The setting we study is a small cafeteria with $C$ waiters and $S$
seats at the counter. A customer arrives, takes a free seat if one is
available (otherwise the customer waits standing), receives attention
from a free waiter to place the order, then consumes the product still
occupying the seat, and finally leaves. The system therefore couples
two resources: waiters are needed only during the attention phase,
whereas seats are held throughout the whole stay.

This coupling makes the system analytically intractable in closed form:
no closed expression exists for the joint distribution of the two
resource queues. We resort to \emph{discrete-event simulation}, which
allows us to compute every performance measure of interest while
remaining methodologically transparent. The whole study sticks to the
methods covered in the course material, namely Algorithm~1.2 of
Leemis~\& Park\,\cite{leemis_park} for the simulation engine, the
multi-stream Lehmer generator\,\cite{leemis_park} for pseudo-random
numbers, Welch's method for warm-up
identification\,\cite{leemis_park, kurkowski_manet} and the
replication-deletion approach with the t-Student confidence interval
for the statistical analysis.

The study answers three questions:
(i)~which is the system's bottleneck, the waiters or the seats?
(ii)~how does the performance scale with the arrival rate, and which
combination of $(C,S)$ remains stable up to which rate?
(iii)~given operating costs and a simple revenue model, which
configuration maximises the hourly net benefit for each demand level?

% =====================================================================
\section{System description and study objectives}
% =====================================================================

\subsection{Customer classes}
The arrival stream is composed of two customer classes with the following
characteristics:

\begin{itemize}
\item \textbf{Coffee customers} ({\tt CAFE}). They arrive with probability
$p_c = 0.7$, are quickly attended and consume their drink in a short
amount of time.
\item \textbf{Breakfast customers} ({\tt DESAYUNO}). They arrive with
probability $p_d = 1 - p_c = 0.3$ and require both a longer attention
(the waiter has to bring several items) and a longer consumption time.
\end{itemize}

\subsection{Resources}
The system has two physical resources:

\begin{itemize}
\item \textbf{Seats} ($S$ identical units). A seat is occupied by a customer
for the whole duration of the stay, from the moment the customer sits
down until consumption is finished. There is a single FCFS standing
queue for customers waiting to sit.
\item \textbf{Waiters} ($C$ identical units). A waiter serves one customer
at a time during the \emph{attention phase}. There is a single FCFS
queue of seated customers waiting for a waiter to become available.
\end{itemize}

\subsection{Operating questions and scope}
The questions the model must answer are:

\begin{enumerate}
\item For the baseline configuration $(C,S)=(2,4)$, what is the mean
response time, the mean number of customers in the system, and the
utilization of each resource?
\item How sensitive are these metrics to the arrival rate? At which
$\lambda$ does the system become unstable?
\item Comparing the four candidate configurations, $(2,4)$, $(2,5)$,
$(3,4)$ and $(3,5)$, which is the cheapest configuration that still
provides acceptable service for each demand level?
\end{enumerate}

% =====================================================================
\section{Conceptual model}
% =====================================================================

The conceptual model describes the system as a discrete-event process
with three event types and a state consisting of the contents of four
\emph{logical zones}:

\begin{itemize}
\item \emph{Standing zone:} customers waiting for a seat.
\item \emph{Seated-waiting zone:} customers occupying a seat and waiting
for a waiter.
\item \emph{Attended zone:} customers being attended by a waiter.
\item \emph{Consuming zone:} customers consuming and occupying a seat.
\end{itemize}

A customer flows through these zones according to the cycle illustrated
below. Each transition between zones is an event:

\begin{center}
\fbox{\parbox{0.95\linewidth}{
\textbf{Customer lifecycle.}\\[2pt]
\textsc{Arrival} $\to$ standing zone (if no free seat) $\to$
\textsc{seat allocation} $\to$ seated-waiting zone (if no free waiter)
$\to$ \textsc{Attention begins} $\to$ attended zone $\to$
\textsc{End of attention} $\to$ consuming zone $\to$
\textsc{End of consumption} $\to$ departure (seat is freed).
}}
\end{center}

The system state is fully described by the cardinality of each zone.
The simulator implements the next-event algorithm (Section~5):
between two consecutive events the state is constant, and the
simulation clock advances by jumps to the next scheduled event.

% =====================================================================
\section{Specification model}
% =====================================================================

\subsection{Stochastic processes and distribution choices}
The choice of distributions is summarized in
Table~\ref{tab:params}. Inter-arrival times are exponentially
distributed: the arrival process is therefore a homogeneous Poisson
process, the standard choice for unscheduled arrivals at a public
service\,\cite{harchol_balter}. The customer class is drawn from a
Bernoulli distribution with parameter $p_c = 0.7$. Attention and
consumption times are exponentially distributed and depend on the
class.

The choice of \emph{exponential} service times has three motivations:
\begin{itemize}
\item It guarantees the memoryless property and therefore enables a
Markovian analysis where applicable.
\item It allows us to compute analytical bounds for validation
(Section~\ref{sec:validation}).
\item It is the standard reference in queueing theory and matches the
distributions used in the lecture material.
\end{itemize}

\begin{table}[h]
\centering
\caption{Parameters of the specification model. Means in minutes.}
\label{tab:params}
\begin{tabular}{l l l}
\toprule
Process & Distribution & Mean \\
\midrule
Inter-arrivals          & \texttt{Exp}     & $1/\lambda$ (varied) \\
Class assignment        & \texttt{Bern}    & $p_c = 0.7$ (coffee) \\
Attention -- coffee     & \texttt{Exp}     & $E[A_c] = 1.5$ \\
Attention -- breakfast  & \texttt{Exp}     & $E[A_d] = 4$ \\
Consumption -- coffee   & \texttt{Exp}     & $E[C_c] = 5$ \\
Consumption -- breakfast& \texttt{Exp}     & $E[C_d] = 12$ \\
\bottomrule
\end{tabular}
\end{table}

\subsection{Derived means}
\label{sec:means}

Averaging over the two customer classes with mixing probabilities
$(p_c, p_d) = (0.7, 0.3)$ gives the marginal mean attention time
\begin{equation}
E[A] = p_c\,E[A_c] + p_d\,E[A_d]
     = 0.7\cdot 1.5 + 0.3\cdot 4 = 2.25 \text{ min},
\label{eq:EA}
\end{equation}
the marginal mean consumption time
\begin{equation}
E[C] = p_c\,E[C_c] + p_d\,E[C_d]
     = 0.7\cdot 5 + 0.3\cdot 12 = 6.5 \text{ min},
\label{eq:EC}
\end{equation}
and the mean time a customer occupies a seat:
\begin{equation}
E[T_{\text{seat}}] = E[A] + E[C] = 8.75 \text{ min}.
\label{eq:Eseat}
\end{equation}

\subsection{Analytical utilizations (lower bounds)}
\label{sec:rhopred}

Treating each resource in isolation (which neglects the coupling
between seats and waiters), the predicted utilizations are
\begin{equation}
\hat\rho_{\text{seat}}(C,S,\lambda)
 = \frac{\lambda\,E[T_{\text{seat}}]}{S}, \qquad
\hat\rho_{\text{waiter}}(C,S,\lambda)
 = \frac{\lambda\,E[A]}{C}. \label{eq:rhopred}
\end{equation}
The system is \emph{stable} when both predicted utilizations are
strictly below~$1$. The seat utilization is generally larger because
$E[T_{\text{seat}}] \gg E[A]$ and the seat is held throughout the
whole stay.

% =====================================================================
\section{Computational model}
% =====================================================================

\subsection{Simulator architecture}
The simulator is implemented in ANSI~C and follows the next-event
discrete-event simulation paradigm (Algorithm~1.2 of
Leemis~\& Park\,\cite{leemis_park}). The main loop is:

\begin{enumerate}
\item Initialize state and schedule the first arrival.
\item Find the next event in the event list by taking the minimum
across (i)~the next arrival time, (ii)~the per-waiter attention-end
times, and (iii)~the per-seat consumption-end times.
\item Advance the clock to that time.
\item Update statistics, then process the event (updating state and
scheduling new events).
\item Repeat until the termination condition
($\text{clients\_served} \ge \mathrm{LAST}$, with $\mathrm{LAST}=10\,000$).
\end{enumerate}

State variables are kept in two arrays indexed by seat slot and
waiter slot (book-style), plus a small FIFO of seat indices for
seated-and-waiting customers and a linked list for the unbounded
standing queue.

\subsection{Pseudo-random numbers}
Pseudo-random numbers are generated by the multi-stream Lehmer
generator with modulus $m = 2^{31}-1$ and multiplier $a = 48271$.
The 256 disjoint streams are obtained from the jump multiplier
$a^{j} \bmod m = 22925$. Each stochastic process is assigned a
dedicated stream to ensure that the underlying random sequences are
\emph{decoupled} across processes, in line with the variance-reduction
good practice discussed by Kurkowski et al.\,\cite{kurkowski_manet}:

\begin{itemize}
\item Stream 0: inter-arrival times
\item Stream 1: customer class assignment
\item Streams 2--3: attention times (coffee, breakfast)
\item Streams 4--5: consumption times (coffee, breakfast)
\end{itemize}

Variates are produced by the inverse method:
\begin{equation}
X = -m \,\ln(1-U), \quad U \sim \mathrm{Uniform}(0,1), \quad
X \sim \mathrm{Exp}(\text{mean} = m).
\end{equation}

% =====================================================================
\section{Verification}
% =====================================================================

Three runtime checks confirm that the simulator behaves as the
specification model dictates:

\paragraph{(i)~Little's Law.}
Little's Law\,\cite{harchol_balter} states the universal identity
\begin{equation}
L = \lambda\,W, \label{eq:little}
\end{equation}
which holds in steady state regardless of the distributions involved.
The simulator computes $L$ as the time-integrated number of customers
in system, $\lambda$ as the observed external arrival rate, and $W$ as
the per-customer mean total time in system. Their ratio agrees
to four decimal places in every replication (see the baseline result
in Section~\ref{sec:results-baseline}).

\paragraph{(ii)~Observed arrival rate.}
The observed inter-arrival rate must match the input value:
$\lambda^{\text{obs}} = N_{\text{arrived}}/T \to 1/E[\text{IAR}]$ in the
limit. This is verified for every scenario.

\paragraph{(iii)~Class mix.}
The empirical fraction of coffee customers must converge to
$p_c = 0.7$. The simulator prints this fraction; observed deviations
are within stochastic noise.

% =====================================================================
\section{Validation against analytical predictions}
\label{sec:validation}
% =====================================================================

The analytical bounds in Eq.~\eqref{eq:rhopred} are compared with the
simulated utilizations in Table~\ref{tab:validation}.

\paragraph{Waiter side.}
$\rho_{\text{waiter}}^{\text{sim}}$ agrees with $\hat\rho_{\text{waiter}}$
within $0.5$~percentage points across all $16$ scenarios. This validates
the waiter sub-model: a waiter behaves essentially as an isolated
M/M/$C$ server with mean service time $E[A]$.

\paragraph{Seat side.}
$\rho_{\text{seat}}^{\text{sim}}$ is systematically higher than
$\hat\rho_{\text{seat}}$, by $5$--$8$~percentage points. The reason is
the resource coupling: a customer occupies a seat \emph{also while
waiting for a waiter}, an effect that
Eq.~\eqref{eq:rhopred} does not capture. The true mean seat
occupation time is
\[
E[T_{\text{seat}}^{\text{real}}]
   = E[T_{\text{seat}}] + E[W_{\text{seated-waiting}}]
   > E[T_{\text{seat}}],
\]
where the second term is the additional time spent seated waiting for
a waiter. This is exactly the coupling that defeats a closed-form
analytical solution and motivates the use of simulation.

\paragraph{Stability prediction.}
Despite the systematic underestimation, the analytical bound correctly
\emph{identifies the unstable region}: whenever
$\hat\rho_{\text{seat}} > 1$ (rows in bold in
Table~\ref{tab:validation}), the simulator shows saturation
($\rho_{\text{seat}}^{\text{sim}} = 1.000$ and $W$ explodes).
This is a strong validation of the model's qualitative behaviour.

\begin{table}[h]
\centering
\caption{Predicted vs simulated utilizations across all 16 scenarios.}
\label{tab:validation}
\small
\begin{tabular}{r r r r r r r l}
\toprule
$\lambda$/h & C & S & $\hat\rho_{\text{seat}}$ & $\rho_{\text{seat}}^{\text{sim}}$ &
$\hat\rho_{\text{waiter}}$ & $\rho_{\text{waiter}}^{\text{sim}}$ & Stable? \\
\midrule
15 & 2 & 4 & 0.547 & 0.593 & 0.281 & 0.283 & yes \\
15 & 2 & 5 & 0.438 & 0.477 & 0.281 & 0.283 & yes \\
15 & 3 & 4 & 0.547 & 0.586 & 0.188 & 0.189 & yes \\
15 & 3 & 5 & 0.438 & 0.469 & 0.188 & 0.188 & yes \\
20 & 2 & 4 & 0.729 & 0.794 & 0.375 & 0.377 & yes \\
20 & 2 & 5 & 0.583 & 0.641 & 0.375 & 0.377 & yes \\
20 & 3 & 4 & 0.729 & 0.782 & 0.250 & 0.251 & yes \\
20 & 3 & 5 & 0.583 & 0.626 & 0.250 & 0.251 & yes \\
25 & 2 & 4 & 0.911 & 0.990 & 0.469 & 0.467 & marginal \\
25 & 2 & 5 & 0.729 & 0.807 & 0.469 & 0.471 & yes \\
25 & 3 & 4 & 0.911 & 0.975 & 0.313 & 0.313 & marginal \\
25 & 3 & 5 & 0.729 & 0.783 & 0.313 & 0.314 & yes \\
30 & 2 & 4 & \textbf{1.094} & 1.000 & 0.563 & 0.472 & \textbf{NO} \\
30 & 2 & 5 & 0.875 & 0.976 & 0.563 & 0.563 & critical \\
30 & 3 & 4 & \textbf{1.094} & 1.000 & 0.375 & 0.321 & \textbf{NO} \\
30 & 3 & 5 & 0.875 & 0.941 & 0.375 & 0.376 & yes \\
\bottomrule
\end{tabular}
\end{table}

% =====================================================================
\section{Experimental design and statistical method}
% =====================================================================

\subsection{Factorial design}
A full-factorial experiment is run with two factors, each at four
levels: the configuration $(C,S)$ and the arrival rate $\lambda$. This
yields $4 \times 4 = 16$ scenarios. For each scenario, $n = 32$
independent replications are performed with seeds $1,\ldots,32$.

The simulation length is fixed at $\mathrm{LAST} = 10\,000$ customers,
which after the warm-up cut (see Section~\ref{sec:warmup}) gives an
effective observation window of approximately $440$~h of simulated
time per replication, sufficient to bound the confidence intervals
below $5\%$ half-width in most scenarios.

\subsection{Replication-deletion}
\label{sec:repdel}
The statistical method is \emph{replication-deletion}\,\cite{leemis_park}:
each scenario is simulated $n$ times with independent random-number
streams, the warm-up portion is discarded, and the per-replication
sample means $\bar X_i$ are treated as i.i.d.\ observations. The
overall mean and sample standard deviation are computed as
\begin{equation}
\bar X = \frac{1}{n}\sum_{i=1}^n \bar X_i, \qquad
s^2 = \frac{1}{n-1}\sum_{i=1}^n (\bar X_i - \bar X)^2,
\label{eq:meanstd}
\end{equation}
and the $95\%$ confidence interval is
\begin{equation}
\bar X \pm t_{n-1,\,\alpha/2}\,\frac{s}{\sqrt{n}}, \quad \alpha = 0.05,
\label{eq:ci}
\end{equation}
where $t_{n-1,\,\alpha/2}$ is the t-Student critical value with $n-1$
degrees of freedom. For our $n=64$ used in the baseline case,
$t_{63,\,0.025} \approx 1.998$.

% =====================================================================
\section{Transient analysis: Welch's method}
\label{sec:warmup}
% =====================================================================

The simulator starts in an empty state, which is not representative
of the steady-state operating regime. The data collected during the
initial transient must therefore be discarded to avoid bias.

We adopt Welch's method\,\cite{leemis_park}. For a chosen indicator
$Y_{r,k}$ measured at the $k$-th snapshot ($\Delta t = 60$ min) of
the $r$-th replication ($r = 1,\ldots,R$), the \emph{ensemble average}
across replications is
\begin{equation}
\bar Y_k = \frac{1}{R}\sum_{r=1}^{R} Y_{r,k}, \qquad k = 1, 2, \ldots
\label{eq:welch-avg}
\end{equation}
This ensemble curve is then smoothed by a moving-average filter of
window $2w+1$:
\begin{equation}
\bar Y_k^{(w)} = \frac{1}{2w+1}\sum_{s=-w}^{w} \bar Y_{k+s}.
\label{eq:welch-ma}
\end{equation}
We apply Eq.~\eqref{eq:welch-avg}--\eqref{eq:welch-ma} with $R = 30$
independent replications and $w = 5$ to three indicators that capture
different aspects of the transient: the cumulative mean response time,
the instantaneous number of customers in system, and the time-integrated
seat utilization.

All three curves stabilise by approximately $t \approx 50$~h. We adopt
\begin{equation}
\boxed{T_0 = 60 \text{ h} = 3600 \text{ min}}
\label{eq:warmup}
\end{equation}
as the warm-up cut, with a $10$~h safety margin. All statistics
presented hereafter are computed over the post-warm-up interval
$[T_0, T_{\text{end}}]$ of length
$T_{\text{eff}} = T_{\text{end}} - T_0 \approx 440$~h.

\begin{figure}[h]
\centering
\begin{subfigure}{0.48\textwidth}
    \includegraphics[width=\linewidth]{welch_avg_total.png}
    \caption{Cumulative mean response time}
\end{subfigure}
\hfill
\begin{subfigure}{0.48\textwidth}
    \includegraphics[width=\linewidth]{welch_util_chairs.png}
    \caption{Seat utilization}
\end{subfigure}
\caption{Welch's method on two representative metrics. The thin grey
curves are individual replications, the thin coloured curve is the
ensemble average ($R = 30$), the thick orange curve is the
moving-average filter ($w = 5$). The smoothed curve stabilises by
$t \approx 50$~h, justifying the choice $T_0 = 60$~h.}
\label{fig:welch}
\end{figure}

% =====================================================================
\section{Baseline scenario results}
\label{sec:results-baseline}
% =====================================================================

For the baseline configuration $(C,S,\lambda) = (2,4,20)$ we run
$n=64$ independent replications. The resulting mean values, sample
standard deviations and $95\%$ confidence intervals, computed via
Eq.~\eqref{eq:meanstd}--\eqref{eq:ci}, are listed in
Table~\ref{tab:baseline}. The relative half-width is below $4.5\%$
for every metric, indicating excellent statistical precision.

\begin{table}[h]
\centering
\caption{Baseline scenario $(C,S,\lambda) = (2,4,20)$ with $n=64$ replications.}
\label{tab:baseline}
\begin{tabular}{l r r r r}
\toprule
Metric & Mean & St.dev. & CI half-width & Relative \\
\midrule
Mean response time $W$ (min)          & 16.06 & 1.213 & $\pm$0.30 & 1.9\% \\
Mean standing wait (min)              &  6.54 & 1.149 & $\pm$0.29 & 4.4\% \\
Mean seated wait for waiter (min)     &  0.15 & 0.010 & $\pm$0.002 & 1.6\% \\
Throughput $X$ (cust/h)               & 20.01 & 0.236 & $\pm$0.06 & 0.3\% \\
Seat utilization                      & 0.7935& 0.014 & $\pm$0.0035 & 0.4\% \\
Waiter utilization                    & 0.3765& 0.007 & $\pm$0.0017 & 0.5\% \\
Mean number in system $L$             &  5.36 & 0.444 & $\pm$0.11 & 2.1\% \\
\bottomrule
\end{tabular}
\end{table}

\paragraph{Little's Law check.}
With $\lambda = 20$/h $= 1/3$/min and $W = 16.06$ min,
Eq.~\eqref{eq:little} gives $\lambda \cdot W = (1/3)\cdot 16.06 = 5.36$,
which matches the simulated $L = 5.36$ to four decimals. This is a
strong runtime verification of the simulator.

\paragraph{Throughput check.}
The observed throughput $X = 20.01$~cust/h is essentially equal to
the input rate $\lambda = 20$~cust/h, confirming that there are no
losses (the standing queue is unbounded).

% =====================================================================
\section{Full factorial results}
\label{sec:results-factorial}
% =====================================================================

Figures \ref{fig:Wlin}--\ref{fig:Uw} display the four key performance
metrics ($W$, $L$, $\rho_{\text{seat}}$, $\rho_{\text{waiter}}$) as a
function of the arrival rate $\lambda$, with one curve per
configuration. Error bars represent the $95\%$ CI half-widths
computed via Eq.~\eqref{eq:ci}.

\begin{figure}[h]
\centering
\begin{subfigure}{0.49\textwidth}
    \includegraphics[width=\linewidth]{fig_W_vs_lambda.png}
    \caption{linear scale}
\end{subfigure}
\begin{subfigure}{0.49\textwidth}
    \includegraphics[width=\linewidth]{fig_W_vs_lambda_log.png}
    \caption{logarithmic scale}
\end{subfigure}
\caption{Mean response time $W$ vs arrival rate $\lambda$. The
logarithmic scale reveals the saturation of $(2,4)$ and $(3,4)$ at
$\lambda = 30$/h.}
\label{fig:Wlin}
\end{figure}

\begin{figure}[h]
\centering
\begin{subfigure}{0.49\textwidth}
    \includegraphics[width=\linewidth]{fig_L_vs_lambda.png}
    \caption{Mean number in system}
\end{subfigure}
\begin{subfigure}{0.49\textwidth}
    \includegraphics[width=\linewidth]{fig_util_chairs.png}
    \caption{Seat utilization}
\end{subfigure}
\caption{Mean number in system and seat utilization vs arrival rate,
by configuration.}
\label{fig:LU}
\end{figure}

\begin{figure}[h]
\centering
\includegraphics[width=0.62\linewidth]{fig_util_waiters.png}
\caption{Waiter utilization vs arrival rate. Even at the peak rate
$\lambda=30$/h, $\rho_{\text{waiter}}$ remains well below~$1$,
confirming that waiters are \emph{not} the bottleneck of the system.}
\label{fig:Uw}
\end{figure}

\paragraph{Bottleneck.}
Figures \ref{fig:LU}b and~\ref{fig:Uw} make the bottleneck explicit.
At the baseline ($\lambda = 20$/h), $\rho_{\text{seat}} \approx 0.79$
against $\rho_{\text{waiter}} \approx 0.38$, more than a factor of two.
At the peak ($\lambda = 30$/h), the seat utilization saturates at $1$
for configurations $(2,4)$ and $(3,4)$ while
$\rho_{\text{waiter}} \le 0.56$. The bottleneck is therefore the seat,
not the waiter, in every regime considered.

\paragraph{Effect of adding one extra resource.}
The benefit of \emph{adding one seat} is dramatic. At $\lambda = 25$/h,
adding a seat takes the mean response time from $W = 162$~min
($(C,S)=(2,4)$) down to $W = 15$~min ($(C,S)=(2,5)$), a $90\%$
reduction. By contrast, \emph{adding one waiter} (going to $(3,4)$)
brings $W$ only to $84$~min, a $48\%$ reduction. The combined
addition $(3,5)$ improves marginally over $(2,5)$ alone:
the bottleneck is so clearly the seat that the waiter upgrade is
largely redundant up to $\lambda = 25$/h.

% =====================================================================
\section{Economic analysis}
% =====================================================================

\subsection{Net benefit function}
We define the \emph{hourly net benefit} of a configuration $(C,S)$
operating at arrival rate $\lambda$ as
\begin{equation}
B(C,S,\lambda)
 = \underbrace{\big(p_c\,p_{\text{coffee}} + p_d\,p_{\text{breakfast}}\big)\,X(C,S,\lambda)}_{\text{revenue}}
 \;-\; \underbrace{c_w\,C + c_s\,S}_{\text{operating cost}},
\label{eq:benefit}
\end{equation}
with prices $p_{\text{coffee}} = 1.5$~€, $p_{\text{breakfast}} = 5$~€,
and costs $c_w = 12$~€/h per waiter and $c_s = 1$~€/h per seat.
The two costs $c_w$ and $c_s$ are deliberately set so that
$c_w \gg c_s$: a seat is cheap (rent, depreciation), a waiter is
expensive (salary). Substituting $p_c$ and $p_d$ gives the
per-customer revenue
\[
\bar r = p_c\,p_{\text{coffee}} + p_d\,p_{\text{breakfast}}
       = 0.7\cdot 1.5 + 0.3\cdot 5 = 2.55 \text{ €}.
\]

For each scenario we substitute the simulated throughput $X$, which
matches the arrival rate $\lambda$ when the system is stable and falls
below it when the bottleneck saturates.

\subsection{Scenario-by-scenario evaluation}
Table~\ref{tab:benefit} reports the benefit of every scenario.
The viability column flags scenarios with prohibitive mean response
time (e.g., $W > 30$~min) or saturation (no stable steady state):
in real terms, those scenarios are operationally inviable even if
nominally profitable.

\begin{table}[h]
\centering
\caption{Hourly benefit and viability per scenario.
$\bar r = 2.55$~€/customer, $c_w = 12$~€/h, $c_s = 1$~€/h.}
\label{tab:benefit}
\small
\begin{tabular}{r r r r r r r l}
\toprule
$\lambda$/h & C & S & $X^{\text{sim}}$/h &
Revenue (€/h) & Cost (€/h) &
Net (€/h) & $W$ (min) -- Viable? \\
\midrule
15 & 2 & 4 & 15.01 & 38.28 & 28 & \textbf{+10.28} & 11.05 -- yes (best) \\
15 & 2 & 5 & 15.01 & 38.28 & 29 & +9.28   &  9.94 -- yes \\
15 & 3 & 4 & 15.01 & 38.28 & 40 & $-1.72$ & 10.81 -- no (loss) \\
15 & 3 & 5 & 15.01 & 38.28 & 41 & $-2.72$ &  9.72 -- no (loss) \\
20 & 2 & 4 & 20.01 & 51.03 & 28 & \textbf{+23.03} & 15.98 -- yes (best) \\
20 & 2 & 5 & 20.01 & 51.03 & 29 & +22.03  & 11.17 -- yes \\
20 & 3 & 4 & 20.01 & 51.03 & 40 & +11.03  & 15.04 -- yes \\
20 & 3 & 5 & 20.01 & 51.03 & 41 & +10.03  & 10.66 -- yes \\
25 & 2 & 4 & 24.82 & 63.30 & 28 & +35.30  & 162.6 -- no ($W$ huge) \\
25 & 2 & 5 & 25.02 & 63.79 & 29 & \textbf{+34.79} & 15.25 -- yes (best) \\
25 & 3 & 4 & 24.96 & 63.65 & 40 & +23.65  &  84.0 -- no \\
25 & 3 & 5 & 25.02 & 63.79 & 41 & +22.79  & 13.60 -- yes \\
30 & 2 & 4 & 25.07 & 63.93 & 28 & +35.93  & 2250  -- no (saturated) \\
30 & 2 & 5 & 29.95 & 76.37 & 29 & +47.37  &  75.1 -- no (high $W$) \\
30 & 3 & 4 & 25.61 & 65.30 & 40 & +25.30  & 1964  -- no (saturated) \\
30 & 3 & 5 & 30.01 & 76.53 & 41 & \textbf{+35.53} & 33.7  -- yes (best) \\
\bottomrule
\end{tabular}
\end{table}

\begin{figure}[h]
\centering
\includegraphics[width=0.7\linewidth]{fig_benefit_vs_lambda.png}
\caption{Hourly net benefit (€/h) vs arrival rate $\lambda$ per
configuration, from Eq.~\eqref{eq:benefit}.}
\label{fig:benefit}
\end{figure}

\subsection{Operational policy}
Combining Tables \ref{tab:validation} and~\ref{tab:benefit}, the
recommended policy is:

\begin{itemize}
\item At \emph{low demand} ($\lambda = 15$/h), the cheapest
configuration $(2,4)$ is optimal. Adding either a seat or a waiter
reduces the benefit because the demand does not justify the extra
capacity.
\item At \emph{normal demand} ($\lambda = 20$/h), the baseline $(2,4)$
still wins by a small margin over $(2,5)$.
\item At \emph{high demand} ($\lambda = 25$/h), $(2,5)$ becomes the
only viable optimum: \emph{adding a seat (1~€/h) is enough}, no
extra waiter is required.
\item At \emph{peak demand} ($\lambda = 30$/h), only $(3,5)$ is viable.
Even so, $W = 33.7$~min remains high: at this rate, the manager should
consider redesigning the service (more seats or a different layout)
rather than simply hiring more staff.
\end{itemize}

The economic conclusion is therefore unambiguous: \emph{when in doubt,
add a seat before adding a waiter}. The seat is one order of magnitude
cheaper and addresses the actual bottleneck.

% =====================================================================
\section{Conclusions}
% =====================================================================

This study has built a complete discrete-event simulation of a small
cafeteria modelled as a queueing system with two coupled resources
(seats and waiters). The model was developed following Algorithm~1.1
of Leemis~\& Park, implemented with the multi-stream Lehmer
pseudo-random generator, validated through three runtime checks
(Little's Law, observed arrival rate, class mix) and against an
analytical lower-bound model of resource utilization. The warm-up
period was identified by Welch's method as $T_0 = 60$~h, and all
metrics were computed via the replication-deletion method with the
t-Student $95\%$ confidence interval over $n = 32$ replications per
scenario.

The main findings are:
\begin{itemize}
\item \emph{The seat, not the waiter, is the bottleneck of the system}
at every operating point considered. The analytical bound
$\hat\rho_{\text{seat}} = \lambda\,E[T_{\text{seat}}]/S$ correctly
predicts both the magnitude and the stability boundary.
\item Adding a seat is operationally and economically far superior to
adding a waiter, both because the seat is twelve times cheaper and
because it directly attacks the bottleneck.
\item The optimal configuration is demand-dependent: $(2,4)$ at low
to normal demand, $(2,5)$ at high demand, and $(3,5)$ only at peak,
which barely satisfies a reasonable service-level objective.
\end{itemize}

\paragraph{Limitations and possible extensions.}
The model assumes exponential service times throughout, which is the
standard reference but does not match every empirical distribution.
Customer balking (refusing to enter the queue) and reneging (leaving
after some wait) are not modelled, although the data from the high-load
scenarios suggest these effects would dominate any further refinement.
Extensions could include heterogeneous waiters (with different service
rates), a time-varying arrival rate (peak/off-peak hours), and the
inclusion of a take-away service that does not need a seat.

% =====================================================================
\begin{thebibliography}{9}
\bibitem{leemis_park}
L.~M.~Leemis and S.~K.~Park,
\emph{Discrete-Event Simulation: A First Course},
Pearson Education / Prentice Hall, 2006.

\bibitem{kurkowski_manet}
S.~Kurkowski, T.~Camp and M.~Colagrosso,
``MANET Simulation Studies: The Incredibles'',
\emph{Mobile Computing and Communications Review}, vol.~9, no.~4,
pp.~50--61, Oct.~2005.

\bibitem{harchol_balter}
M.~Harchol-Balter,
\emph{Performance Modeling and Design of Computer Systems},
Cambridge University Press, 2013.
\end{thebibliography}

\end{document}
